\renewcommand{\theequation}{A.\arabic{section}.\arabic{equation}}
\renewcommand{\thesection}{A-\arabic{section}}
\setcounter{equation}{0}
\setcounter{section}{0}
\section*{Appendix}

\section{Household First-Order Conditions}
\label{AppHouseholdFOCs}

This appendix derives the household first-order conditions of the model presented in Section~\ref{sec:model}. Zubairy's original derivations, which apply unchanged to the parts of the model that are not modified by our tax extensions, are reproduced in \cn{Zubairy2010}. We restrict attention here to the conditions that are either new or modified relative to the baseline.

\subsection{Utility and Constraints}

The representative household maximizes:
\begin{equation}\label{AppEqUtility}
    E_0 \sum_{t=0}^{\infty} \beta^t\, d_t\, U(x_t^c, h_t) \qquad
    \text{with}\qquad U(x_t^c, h_t) = \frac{1}{1-\gamma}\left[(x_t^c)^a (1-h_t)^{1-a}\right]^{1-\gamma}
\end{equation}
where $x_t^c = c_t - b^c s_{t-1}^C$ is habit-adjusted consumption, $s_t^C = \theta^c s_{t-1}^C + (1-\theta^c) c_t$ is the private habit stock, $a$ is pinned down in steady state, $\gamma$ is risk aversion, and $b^c \in (0,1)$ is the deep-habit parameter.

The household faces the flow budget constraint:
\begin{equation}\label{AppEqBC}
\begin{split}
    (1+\tau_t^c)c_t + \frac{b_t}{\pi_t} + (1 - \tau_t^{ic} - \tau_t^k e_{\tau,t})i_t
    = & \; (1-\tau_t^w)\frac{w_t}{\tilde w_t}h_t + (1-\tau_t^k)r_t^k u_t k_{t-1} + \tau_t^k \delta_\tau k_{t-1}^\tau \\
    & + \frac{[1 + (R_{t-1}-1)(1-\tau_t^i)]\, b_{t-1}}{\pi_t} + (1-\tau_t^d)div_t + tr_t
\end{split}
\end{equation}
the capital accumulation law:
\begin{equation}\label{AppEqCapK}
    k_t = (1-\delta)k_{t-1} + i_t\left[1 - \tfrac{\kappa}{2}\left(\tfrac{\mu_t i_t}{i_{t-1}} - 1\right)^2\right]
\end{equation}
and the tax-basis accumulation law (new):
\begin{equation}\label{AppEqTaxBasis}
    k_t^\tau = (1 - \delta_\tau) k_{t-1}^\tau + i_t (1 - e_{\tau,t})
\end{equation}
The Lagrange multipliers on \eqref{AppEqBC}, \eqref{AppEqCapK}, and \eqref{AppEqTaxBasis} are $\lambda_t$, $\lambda_t q_t$, and $\lambda_t v_{\tau,t}$ respectively. In Zubairy's baseline the tax basis of capital is not separated from the physical capital stock; here we introduce it because immediate expensing ($e_\tau$) drives a wedge between the two.

\subsection{Marginal Utility of Consumption}

The FOC with respect to $c_t$ gives $\partial U/\partial c_t = (1+\tau_t^c)\lambda_t$. Substituting the marginal utility from \eqref{AppEqUtility}:
\begin{equation}\label{AppEqLambda}
    \lambda_t = d_t\, a\, \frac{\left[(x_t^c)^a (1-h_t)^{1-a}\right]^{1-\gamma}}{x_t^c (1+\tau_t^c)}
\end{equation}
Relative to \cn{Zubairy2010}, the $(1+\tau_t^c)$ factor in the denominator is new. Because $\lambda_t$ enters every forward-looking equation of the household problem, the consumption tax propagates into the bond Euler, capital Euler, investment FOC, wage Phillips curve, and pricing block without changing the functional form of those equations.

\subsection{Labor--Leisure Condition}

The FOC with respect to $h_t$ combined with \eqref{AppEqLambda} gives the intratemporal condition:
\begin{equation}\label{AppEqFOClabor}
    d_t (1-a) \frac{\left[(x_t^c)^a (1-h_t)^{1-a}\right]^{1-\gamma}}{1 - h_t}
    = \lambda_t \frac{w_t (1-\tau_t^w)}{\tilde w_t}
\end{equation}
This is functionally identical to Zubairy's condition after multiplication by $(1+\tau_t^c)$, which reappears through the modified $\lambda_t$.

\subsection{Bond Euler Equation}

Nominal bonds pay $R_t$ in period $t+1$ per dollar invested, but only the interest earned, $R_t - 1$, is taxable. Combining the FOC with respect to $b_t$ and applying \eqref{AppEqLambda}:
\begin{equation}\label{AppEqBondEuler}
    \lambda_t = \beta \left[1 + (R_t - 1)(1 - \tau_{i,t+1})\right] \frac{\lambda_{t+1}}{\pi_{t+1}}
\end{equation}
Under $\tau_i = 0$ this collapses to Zubairy's Euler equation.

\subsection{Utilization and Capital Euler}

The FOC with respect to utilization $u_t$ delivers:
\begin{equation}\label{AppEqUtil}
    (1-\tau_t^k) r_t^k = \gamma_1 u_t^{\sigma_u}
\end{equation}
The FOC with respect to $k_t$ (using the capital constraint) gives:
\begin{equation}\label{AppEqCapEuler}
    \lambda_t q_t = \beta E_t \lambda_{t+1}
    \left[(1 - \tau_{k,t+1}) r_{t+1}^k u_{t+1}
     - \tfrac{\gamma_1}{1+\sigma_u}(u_{t+1}^{1+\sigma_u} - 1)
     + q_{t+1}(1 - \delta)\right]
\end{equation}
The depreciation deduction $\delta q_{t+1} \tau_{t+1}^k u_{t+1}$ present in \cn{Zubairy2010}, which credits the capital-tax shield directly in the Euler, is absent here. Instead the shield is carried entirely through the shadow value $v_{\tau,t}$ derived below, avoiding double-counting.

\subsection{Shadow Value of the Tax Basis}

The FOC with respect to $k_t^\tau$ delivers a forward-looking recursion. Because an extra unit of tax basis today reduces next-period taxes by $\tau_{k,t+1}\delta_\tau$ and leaves $(1 - \delta_\tau)$ units of remaining basis, we obtain:
\begin{equation}\label{AppEqVTau}
    \lambda_t v_{\tau,t} = \beta E_t
    \left[\lambda_{t+1}\left((1 - \delta_\tau) v_{\tau,t+1} + \tau_{k,t+1}\delta_\tau\right)\right]
\end{equation}
This equation is new relative to \cn{Zubairy2010} and is what makes it possible to separate immediate expensing (through $e_\tau$) from the ordinary depreciation schedule.

\subsection{Investment First-Order Condition}

The FOC with respect to $i_t$ is the most extensively modified condition. The out-of-pocket cost of a unit of investment falls by $\tau_t^{ic}$ (the ITC) and by $\tau_t^k e_{\tau,t}$ (the immediate expensing wedge). The non-expensed portion $(1 - e_{\tau,t})$ generates future depreciation deductions valued at $v_{\tau,t}$:
\begin{equation}\label{AppEqInvFOC}
\begin{split}
    \lambda_t (1 - \tau_t^{ic} - \tau_t^k e_{\tau,t}) = & \; \lambda_t q_t \left[1 - \tfrac{\kappa}{2}\!\left(\tfrac{\mu_t i_t}{i_{t-1}} - 1\right)^2 - \kappa \tfrac{\mu_t i_t}{i_{t-1}}\!\left(\tfrac{\mu_t i_t}{i_{t-1}} - 1\right)\right] \\
    & + \beta E_t \left[\lambda_{t+1} q_{t+1} \kappa \tfrac{\mu_{t+1} i_{t+1}^2}{i_t^2}\!\left(\tfrac{\mu_{t+1} i_{t+1}}{i_t} - 1\right)\right] \\
    & + \lambda_t v_{\tau,t}(1 - e_{\tau,t})
\end{split}
\end{equation}
Under $\tau^{ic} = e_\tau = 0$ the LHS and last term collapse to Zubairy's original expression.

\subsection{Dividend Tax}

The dividend tax $\tau_t^d$ does not appear in any household first-order condition. Dividends are a residual object taken as given by the household:
\begin{equation}\label{AppEqDiv}
    div_t = y_t - w_t h_t - r_t^k u_t k_{t-1}
\end{equation}
The dividend tax only affects the household's after-tax income (equation~\ref{AppEqBC}) and enters the government budget constraint through tax revenue.

\newpage

\section{Steady-State Derivations}
\label{AppSSDerivations}

We compute the deterministic steady state analytically by setting all shocks to their unconditional means ($z = d = \mu = 1$; all fiscal innovations equal to zero) and imposing stationarity. Below, a bar denotes a steady-state value and we drop time subscripts.

\subsection{Shadow Value of the Tax Basis in Steady State}

Setting the recursion \eqref{AppEqVTau} to a fixed point and using $\bar{\lambda}_t = \bar{\lambda}_{t+1} = \bar{\lambda}$:
\begin{equation}\label{AppEqSSvtau}
    \bar{v}_\tau = \frac{\beta \bar{\tau}^k \delta_\tau}{1 - \beta(1 - \delta_\tau)}
\end{equation}
which is the present value of the geometric stream of depreciation deductions on one unit of tax basis. This expression is valid for any positive $\bar{\tau}^k$ and $\delta_\tau$, and does not depend on the immediate expensing rate.

\subsection{Tobin's $q$ in Steady State}

Setting the investment FOC \eqref{AppEqInvFOC} to a fixed point (all adjustment cost terms vanish because $\bar{\mu} = 1$ and $\bar{i}_t = \bar{i}_{t-1}$):
\begin{equation}\label{AppEqSSq}
    \bar{q} = 1 - \bar{\tau}^{ic} - \bar{\tau}^k \bar{e}^\tau - \bar{v}_\tau (1 - \bar{e}^\tau)
\end{equation}
Under $\bar{\tau}^{ic} = 0$ and $\bar{e}^\tau = 0$, this collapses to $\bar{q} = 1 - \bar{v}_\tau$, matching the geometric-stream expression in Zubairy's baseline.

\subsection{Real Return on Capital, Real Interest Rate, and Real Wage}

The steady-state capital Euler \eqref{AppEqCapEuler} combined with $\bar{u} = 1$ implies:
\begin{equation}\label{AppEqSSrk}
    \bar{r}^k = \frac{\bar{q}}{\beta(1 - \tau^k)} - \frac{q(1-\delta)}{1-\tau^k}
\end{equation}
The steady-state bond Euler \eqref{AppEqBondEuler} with $\bar{\pi} = 1$ delivers:
\begin{equation}\label{AppEqSSR}
    \bar{R} = 1 + \frac{1/\beta - 1}{1 - \bar{\tau}^i}
\end{equation}
The real wage follows from the firm's optimization conditions and, given the Cobb--Douglas production function, from the capital--labor ratio and marginal cost:
\begin{equation}\label{AppEqSSw}
    \bar{w} = (1-\theta)\,\overline{mc}\left(\frac{\bar{r}^k}{\theta \overline{mc}}\right)^{-\theta/(1-\theta)}
\end{equation}
where $\overline{mc} = (\eta - 1)/\eta$ is pinned down by the desired steady-state markup.

\subsection{Aggregate Ratios}

Given $\bar{h}$, the steady-state capital stock and output follow:
\begin{equation}\label{AppEqSSKY}
    \frac{\bar{k}}{\bar{h}} = \left(\frac{\theta \overline{mc}}{\bar{r}^k}\right)^{1/(1-\theta)}, \qquad
    \bar{y} = \bar{k}^\theta \bar{h}^{1-\theta} - \psi
\end{equation}
The fixed operating cost $\psi$ is set residually to hit the target dividend share $\overline{div}/\bar{y}$; steady-state investment is $\bar{i} = \delta \bar{k}$; steady-state consumption is $\bar{c} = \bar{y} - \bar{i} - \bar{g}$, using the resource constraint with steady-state utilization cost equal to zero; and steady-state government purchases and debt are determined by the calibrated ratios $\bar{g}/\bar{y}$ and $\bar{b}/\bar{y}$.

\subsection{Marginal Utility in Steady State}

The steady state of the marginal utility equation \eqref{AppEqLambda} pins down the consumption share parameter $a$ jointly with $\bar{\lambda}$. Given $\bar{h}$, $\bar{x}^c = (1 - b^c)\bar{c}$, and the labor--leisure condition \eqref{AppEqFOClabor}, we solve:
\begin{equation}\label{AppEqSSa}
    a = \frac{\bar{x}^c (1 + \bar{\tau}^c)}{\bar{x}^c (1 + \bar{\tau}^c) + (\bar{w}/\tilde{w})(1-\bar{h})}, \qquad
    \bar{\lambda} = a\,\frac{\left[(\bar{x}^c)^a (1-\bar{h})^{1-a}\right]^{1-\gamma}}{\bar{x}^c (1+\bar{\tau}^c)}
\end{equation}
Both quantities depend on estimated parameters and are recomputed at each draw during Bayesian inference.

\subsection{Government Budget in Steady State}

Setting the government budget constraint \eqref{eq:gbc} to a stationary value $\bar{b}$ and using $\bar{\pi} = 1$:
\begin{equation}\label{AppEqSSbudget}
    \bar{b}\left(1 - \bar{R}\right) + \overline{taxrev} = \bar{g} + \bar{tr}
\end{equation}
Given the target debt-to-output ratio and the calibrated tax rates, this determines the steady-state transfer $\bar{tr}$ residually. The debt-residual shock $\varepsilon_t^b$ has mean zero and therefore does not appear in the steady state.
